%% 01-introduction.tex — Why reproducibility matters

\section{Introduction}

Congressional redistricting is one of the few governmental acts that directly
determines who represents whom in Congress for the next decade.
When a research program proposes an algorithmic alternative to human map-drawing,
the burden of proof is correspondingly high: every claim must be independently
verifiable, every result must be exactly reproducible, and every design decision
must be documented so that critics can challenge it and courts can audit it.

This reproducibility package fulfills that obligation.
It provides complete instructions, open-source code, public data sources,
and cryptographic audit trails so that any researcher---or any judge's
technical expert---can reproduce every result in the portfolio from scratch.

\subsection*{Why Redistricting Reproducibility Is Different}

Standard academic reproducibility asks: given the same code and data,
do you get the same numbers?
Redistricting reproducibility asks something stronger: given independent
implementations of the same algorithm applied to Census Bureau data,
do you get byte-identical district assignments on the same hardware architecture?

The answer, for this portfolio, is yes---and this package explains how
to verify it.
The \texttt{redist label-verify} command (Section~\ref{sec:verification})
checks the SHA-256 fingerprint of every input file, every parameter setting,
and every output district assignment, linking them into a tamper-evident chain.

\subsection*{The SHA-256 Audit Chain}

Every result in this portfolio is traceable through a four-step cryptographic chain:

\begin{enumerate}[leftmargin=*, itemsep=4pt]
  \item \textbf{Config hash}: the exact parameter values used
        (population tolerance, METIS seed, edge-weight scaling factor).
  \item \textbf{Build hash}: the SHA-256 of every input Census file,
        confirming no data was altered.
  \item \textbf{Analysis hash}: the SHA-256 of the district assignment output,
        confirming the algorithm ran identically.
  \item \textbf{Report hash}: the SHA-256 of summary statistics and maps,
        confirming downstream analysis was not modified.
\end{enumerate}

\noindent The chain is described in detail in Section~\ref{sec:verification}
and implemented in the \texttt{redist-cli} Rust crate~\cite{dellaLibera2026recursive}.

\subsection*{AEA Compliance}

This package follows the American Economic Association's Data and Code
Availability Policy~\cite{aeaDataPolicy}, which requires:
(1) all data publicly available or provided,
(2) all code posted in a replication archive,
(3) a README describing the steps to reproduce each result,
and (4) expected runtimes on the reference hardware.
Each of these requirements is addressed in the sections below.
